\section{Elektr zaryadlari. Kulon qonuni}



\begin{enumerate}

\item Vakuumda joylashgan nuqtaviy musbat $q_1$ zaryadning o'zidan $r$ masofada joylashgan nuqtaviy musbat $q_2$ zaryadga ko'rsatadigan ta'sir kuchi ifodasini ko'rsating.
\begin{enumerate}[label=\Alph*)]
\item $F = \frac{1}{4\pi\epsilon_0} \frac{q_1 q_2}{r^2}$
\item $F = -k \frac{q_1 q_2}{r^3}$
\item $F = \frac{1}{4\pi\epsilon_0} \frac{q_1 q_2}{r^3}$
\item $F = \frac{1}{4\pi\epsilon_0 \epsilon} \frac{q_1 q_2}{r}$
\end{enumerate}

\item Ikki sharchaga bir xil $2 \cdot 10^{-8} \, C$ zaryad berilgan. Zaryadlar bir-birini $5 \cdot 10^{-4} \, N$ kuch bilan itarishmoqda. Sharchalar orasidagi masofa (sm) qancha?
\begin{enumerate}[label=\Alph*)]
\item 7,5
\item 8,5
\item 8
\item 7
\end{enumerate}

\item Ikkita bir xil sharcha bir-biridan 10 sm masofada turibdi. Ular bir xil miqdorda manfiy zaryadga ega bo'lib, 0,23 mN kuch bilan o'zaro ta'sirlashadi. Har qaysi sharchadagi "ortiqcha" elektronlar sonini toping.
\begin{enumerate}[label=\Alph*)]
\item $1,0 \cdot 10^8$
\item $1,0 \cdot 10^{11}$
\item $1,0 \cdot 10^{14}$
\item $1,0 \cdot 10^{12}$
\end{enumerate}

\item Nuqtaviy zaryadlarning o'zaro ta'sir kuchi 10 mN, zaryad miqdorlari $5,2 \cdot 10^{-6} \, C$ va $5 \cdot 10^{-6} \, C$ bo'lsa, ular orasidagi masofa (m) qanday?
\begin{enumerate}[label=\Alph*)]
\item 10,15
\item 0,31
\item 1
\item 4,83
\end{enumerate}

\item 6,6 mKl va 11 mKl zaryadlar suvda bir-biridan 3,3 sm masofada turib, qanday kuch bilan o'zaro ta'sirlashadi?
\begin{enumerate}[label=\Alph*)]
\item $18 \cdot 10^{-2}$
\item $7,8 \cdot 10^{-2}$
\item $6,4 \cdot 10^{-4}$
\item TJY
\end{enumerate}

\item Bir-biridan 30 sm (markazlari orasidagi) masofada turgan ikkita birday zaryadlangan sharcha o'zaro 10 dina kuch bilan ta'sirlashadi. Sharchalardan biridagi zaryadni toping (pC)? $1 Dina = 10^{-10} \, N$.
\begin{enumerate}[label=\Alph*)]
\item 484
\item 424
\item 252
\item 100
\end{enumerate}

\item Nuqtaviy zaryadlar orasidagi masofa 3 marta o'zgartirilsa, ularning o'zaro ta'sir kuchi necha marta o'zgaradi?
\begin{enumerate}[label=\Alph*)]
\item 3
\item 9
\item 27
\item 81
\end{enumerate}

\item Agar ikki nuqtaviy zaryadning zaryadi to'rt baravar orttirilsa, ular orasidagi masofa esa ikki marta kamaytirilsa, ular orasidagi o'zaro ta'sir kuchi qanday o'zgaradi?
\begin{enumerate}[label=\Alph*)]
\item 8 marta kamayadi
\item 8 marta ortadi
\item 64 marta ortadi
\item 2 marta kamayadi
\end{enumerate}

\item Ikki nuqtaviy zaryad orasidagi masofa 3 marta ortsa, ular orasidagi ta'sir kuchi qanday o'zgaradi?
\begin{enumerate}[label=\Alph*)]
\item 3 marta oshadi
\item 9 marta kamayadi
\item o'zgarmaydi
\item 3 marta kamayadi
\end{enumerate}

\item Vakuumda joylashgan ikkita zaryadlangan zarracha o'zaro ta'sirlashmoqda. Agar zaryadlardan birini ikki marta orttirib, boshqa zaryad 4 marta kamaytirilsa hamda ularning orasidagi masofa 3 marta kamaytirilsa, ular orasidagi Kulon kuchi qanday o'zgaradi?
\begin{enumerate}[label=\Alph*)]
\item 6 marta kamayadi
\item 4,5 marta ortadi
\item 2,5 marta ortadi
\item 9 marta ortadi
\end{enumerate}

\item Radiusi 5 sm bo'lgan va $q = +15 \, nC$ zaryadga ega bo'lgan mis halqa o'z markazida joylashgan $q_1 = -3 \, nC$ nuqtaviy zaryadga qanday kuch bilan ta'sir qiladi?
\begin{enumerate}[label=\Alph*)]
\item 0
\item 40
\item 80
\item 160
\end{enumerate}

\item Havoda 3 sm masofada turgan ikki zaryad bir-biri bilan $2 \cdot 10^{-5} \, N$ kuch orqali ta'sirlashadi. Qandaydir suyuq muhitda shu zaryadlar 5 sm masofadan turib, $1,5 \cdot 10^{-6} \, N$ kuch bilan ta'sirlashadi. Shu suyuq muhitning dielektrik singdiruvchanligini aniqlang.
\begin{enumerate}[label=\Alph*)]
\item $\approx 12$
\item $\approx 5$
\item $\approx 0,2$
\item $\approx 1$
\end{enumerate}

\item Ikkita zaryad suvda bir-biriga 3 sm masofada qanday kuch bilan o'zaro ta'sirlashsa, vakuumda 27 sm masofada ham shunday kuch bilan o'zaro ta'sirlashadi. Suvning dielektrik singdiruvchanligi qancha?
\begin{enumerate}[label=\Alph*)]
\item 81
\item 98
\item 76
\item 9
\end{enumerate}

\item Ikkita zaryad vakuumda $r_1$ masofadan o'zaro ta'sirlashdi. Ularning o'zaro ta'sir kuchini o'zgartirmasdan singdiruvchanligi $\epsilon$ bo'lgan dielektrikda qanday $r_2$ masofada joylashtirish kerak?
\begin{enumerate}[label=\Alph*)]
\item $r_2 = \sqrt{r_1 \epsilon}$
\item $r_2 = \frac{\sqrt{\epsilon}}{r_1}$
\item $r_2 = \frac{r_1}{\sqrt{\epsilon}}$
\item $r_2 = r_1 \sqrt{\epsilon}$
\end{enumerate}

\item Ikkita zaryad havoda bir-biridan $r_1$ masofada turibdi. Ular orasidagi ta'sir kuchi o'zgarmasdan qolishi uchun ularni dielektrik singdiruvchanligi $\epsilon_2$ bo'lgan boshqa muhitda qanday $r_2$ masofada joylashtirish kerak?
\begin{enumerate}[label=\Alph*)]
\item $r_1 = r_2 \sqrt{\epsilon_2}$
\item $r_2 = r_1 \sqrt{\epsilon_2}$
\item $r_2 = 2r_1 \sqrt{\epsilon_2}$
\item $r_1 = r_2 \epsilon_2$
\end{enumerate}

\item $q_1 = q$ va $q_2 = 3q$ zaryadli ikkita bir xil metal sharcha bir-biriga tekkizildi va avvalgi joyiga qaytarildi. Bunda ularning o'zaro ta'sir kuchi necha marta ortadi?
\begin{enumerate}[label=\Alph*)]
\item 9/8
\item 4/3
\item 9/5
\item 9/4
\end{enumerate}

\item $3q$ va $5q$ zaryadli ikkita bir xil metall sharcha bir-biriga tekkizildi va avvalgi vaziyatiga qaytarildi. Bunda ularning o'zaro ta'sir kuchi necha marta ortadi?
\begin{enumerate}[label=\Alph*)]
\item 49/48
\item 16/15
\item 25/21
\item 121/96
\end{enumerate}

\item Biriga N ta, ikkinchisiga N-1 ta elektron yetishmayotgan ikkita bir xil sharcha bir-biridan d masofada turib $F_1$ kuch bilan ta'sirlashyapti. Ular bir-biriga tekkizilib yana oldingi vaziyatiga qaytganda $F_2$ kuch bilan ta'sirlashadi. $F_1/F_2$ nisbat nimaga teng? $N \ge 2$.
\begin{enumerate}[label=\Alph*)]
\item 1
\item $\frac{N(4N-1)}{(2N-1)^2}$
\item $\frac{4N(N-1)}{(2N-1)^2}$
\item $\frac{4N(2N-1)}{(N-1)^2}$
\end{enumerate}

\item Bir xil ishorali $q$ va $4q$ zaryadlar bilan zaryadlangan ikkita bir xil metall sharcha bir-biridan r masofada turibdi. Sharchalar bir-biriga tekkiziladi. O'zaro ta'sir kuchi avvalgidek qolishi uchun ularni qanday masofaga surish lozim?
\begin{enumerate}[label=\Alph*)]
\item 2,5r
\item 0,4r
\item r
\item 1,25r
\end{enumerate}

\item Oralaridagi masofa diametrlaridan juda kichik bo'lgan ikki yassi metall plastina qarama-qarshi ishorali, kattaligi bir xil bo'lgan zaryadlar bilan zaryadlangan va bir-biriga F kuch bilan tortiladi. Agar plastinalar orasidagi masofa 2 marta orttirilsa, bu kuch qanday o'zgaradi?
\begin{enumerate}[label=\Alph*)]
\item 4 marta kamayadi
\item 2 marta kamayadi
\item deyarli o'zgarmaydi
\item $\sqrt{2}$ marta kamayadi
\end{enumerate}

\item Quyidagilarni qaysi biri elektr doimiysini birligi hisoblanadi?
\begin{enumerate}[label=\Alph*)]
\item $kg \cdot m^{-3} \cdot s^{-4} \cdot A^2$
\item $kg \cdot m^2 \cdot s^{-3} \cdot A^3$
\item $kg^{-2} \cdot m^2 \cdot s^3 \cdot A^2$
\item $kg^{-1} \cdot m^{-3} \cdot s^4 \cdot A^2$
\end{enumerate}

\item Ikkita bir xil nuqtaviy $q$ zaryadlar o'rtasiga $q/2$ zaryad joylashtirilsa, $q$ zaryadga ta'sir qiluvchi kuch necha marta ortadi?
\begin{enumerate}[label=\Alph*)]
\item 1,5
\item 2
\item 3
\item 4
\end{enumerate}

\item Uzunligi 100 sm li iplarga osilgan, har birining massasi 0,25 g bo'lgan $q_1 = q_2$ zaryadli ikkita sharcha bir-biridan 6 sm qochgan. Har bir sharchaning zaryadi necha nKl ga teng?
\begin{enumerate}[label=\Alph*)]
\item 5,4
\item 6,7
\item 7,9
\item 10
\end{enumerate}

\item Chizmadagi zaryadlangan sharchalarning o'zaro ta'sir kuchi necha nyuton ekanligini aniqlang? Ulardan biri qo'zg'almas ($g = 10 \, m/s^2, m = 10 \, g, \tan 35^\circ = 0,7$).


\begin{tikzpicture}[scale=1.5, >=Stealth]
    % Ship (Ceiling)
    \fill[pattern=north east lines] (-1,2) rectangle (1,2.1);
    \draw (-1,2) -- (1,2);
    
    % Vertikal osilgan sharcha
    \draw (0,2) -- (0,0.5);
    \draw[fill=white] (0,0.5) circle (0.1) node[below=4pt] {$m$};
    \node at (0,0.5) {$\cdot$};
    
    % Burchak ostidagi sharcha
    \draw (0,2) -- (1,0.5);
    \draw[fill=white] (1,0.5) circle (0.1);
    \node at (1,0.5) {$\cdot$};
    
    % Kuchlar (mg va F_k)
    \draw[->, thick] (1,0.5) -- (1,-0.2) node[right] {$mg$};
    \draw[->, thick] (1,0.5) -- (1.7,0.5) node[above] {$F_k$};
    
    % Burchak 35 gradus
    \draw[dashed] (0,2) -- (0,1);
    \coordinate (top) at (0,2);
    \coordinate (vert) at (0,1);
    \coordinate (side) at (1,0.5);
    \pic [draw, "35$^\circ$", angle eccentricity=1.5, angle radius=0.6cm] {angle = vert--top--side};
\end{tikzpicture}
\begin{enumerate}[label=\Alph*)]
\item 0,07
\item 0,7
\item 7
\item 10
\end{enumerate}

\item Radiusi $5 \, sm$ bo'lgan aylanaga zaryadining moduli $q = 2 \, \mu C$ bo'lgan 3 ta zaryad va markazga zaryadining moduli $Q = 5 \, \mu C$ bo'lgan boshqa zaryad joylashtirildi. Aylananing markazidagi zaryadga ta'sir qiladigan natijaviy Kulon kuchini toping ($\mu N$). ($\alpha = 60^\circ$).

\begin{tikzpicture}[scale=2, >=stealth]
    % Aylana (punktir chiziq bilan)
    \draw[dashed, color=gray] (0,0) circle (1);

    % Markazdagi zaryad (Q+)
    \draw[fill=white] (0,0) circle (0.12) node {$+$};
    \node[left=5pt] at (0,0) {$Q$};

    % Zaryadlarning o'rni (alfa = 60 gradus bo'lganda: 90, 30, -30 graduslar)
    
    % 1. Yuqoridagi musbat zaryad (90 gradusda)
    \draw[dashed] (0,0) -- (90:1);
    \draw[fill=white] (90:1) circle (0.08) node {\tiny $+$};
    \node[above=3pt] at (90:1) {$q$};

    % 2. O'rtadagi manfiy zaryad (30 gradusda)
    \draw[dashed] (0,0) -- (30:1);
    \draw[fill=white] (30:1) circle (0.08) node {\tiny $-$};
    \node[right=3pt] at (30:1) {$q$};

    % 3. Pastdagi musbat zaryad (-30 gradusda)
    \draw[dashed] (0,0) -- (-30:1);
    \draw[fill=white] (-30:1) circle (0.08) node {\tiny $+$};
    \node[below right=2pt] at (-30:1) {$q$};

    % Burchaklarni belgilash (alfa)
    \coordinate (O) at (0,0);
    \coordinate (T) at (90:0.7);
    \coordinate (M) at (30:0.7);
    \coordinate (B) at (-30:0.7);

    \pic [draw, "$\alpha$", angle eccentricity=1.5, angle radius=0.5cm] {angle = M--O--T};
    \pic [draw, "$\alpha$", angle eccentricity=1.5, angle radius=0.5cm] {angle = B--O--M};

\end{tikzpicture}

\begin{enumerate}[label=\Alph*)]
\item 22
\item 32
\item 0
\item 54
\end{enumerate}
\end{enumerate}

\begin{table}[h!]
\centering
\begin{tabular}{|c|c|c|c|c|c|c|c|c|c|}
\hline
№ & J. & № & J. & № & J. & № & J. & № & J. \\ \hline
1 & A & 6 & D & 11 & A & 16 & B & 21 & D \\ \hline
2 & B & 7 & B & 12 & B & 17 & B & 22 & C \\ \hline
3 & B & 8 & C & 13 & A & 18 & C & 23 & B \\ \hline
4 & D & 9 & B & 14 & C & 19 & D & 24 & A\\ \hline
5 & D & 10 & B & 15 & A & 20 & C & 25 &C \\ \hline
\end{tabular}
\end{table}